\subsection{Etap Test}
Automatyzacja testów jest fundamentem zaufania do procesu CI/CD. W projekcie testy są integralną częścią każdego łańcucha budowania i stanowią bramkę blokującą przejście do dalszych, kosztownych etapów.

\subsubsection{Testy integracyjne: JUnit 5 i @QuarkusTest}
Podstawowym narzędziem weryfikacji poprawności logiki biznesowej jest framework \textbf{JUnit 5}, zintegrowany z biblioteką testową Quarkusa poprzez zależność \texttt{quarkus-junit5}.

\textbf{Adnotacja \texttt{@QuarkusTest}:}
W odróżnieniu od klasycznych testów jednostkowych, adnotacja \texttt{@QuarkusTest} uruchamia przed zestawem testów pełny kontekst aplikacji Quarkus -- z kontenerem CDI, warstwą REST i skonfigurowanymi połączeniami do zewnętrznych usług. Testy mogą zatem weryfikować rzeczywiste zachowanie endpointów HTTP przy użyciu zintegrowanej biblioteki \textbf{RestAssured}. Jej API wzoruje się na notacji \textbf{Gherkin} (\textit{Given/When/Then}), stosowanej w metodykach BDD (\textit{Behaviour-Driven Development}): blok \texttt{given()} definiuje warunki wstępne żądania (nagłówki, ciało), \texttt{when()} wykonuje wywołanie HTTP, a \texttt{then()} weryfikuje odpowiedź. Dzięki temu intencja testu jest czytelna bez znajomości samej biblioteki:

\begin{lstlisting}[language=java, caption={Test integracyjny endpointu REST mikroserwisu \textit{users} (apps/users/src/test/.../UserResourceTest.java)}, label={lst:quarkus-test-users}]
@QuarkusTest
class UserResourceTest {

    @Test
    void testListUsers() {
        given()
            .when().get("/users")
            .then()
            .statusCode(200);
    }

    @Test
    void testCreateUser() {
        given()
            .contentType(JSON)
            .body("{\"name\":\"Alice\",\"email\":\"alice@example.com\"}")
            .when().post("/users")
            .then()
            .statusCode(201);
    }
}
\end{lstlisting}

\textbf{Dev Services -- automatyczna infrastruktura testowa:}
Kluczową cechą ekosystemu Quarkusa jest mechanizm \textbf{Dev Services}: jeśli w konfiguracji aplikacji nie podano jawnego adresu URL bazy danych, Quarkus automatycznie uruchamia kontener PostgreSQL przy użyciu biblioteki \textit{Testcontainers}. Programista nie musi konfigurować środowiska testowego ani utrzymywać lokalnej instancji bazy -- wystarczy wywołanie \texttt{./gradlew test}. Testy zawsze działają na izolowanej, świeżej instancji bazy danych, co eliminuje klasyczny problem współdzielonego stanu testowego.

\textbf{Izolacja transakcyjna -- \texttt{@TestTransaction}:}
Dla testów weryfikujących logikę perzystencji stosowana jest adnotacja \texttt{@TestTransaction}: każdy test jest owijany w transakcję bazodanową, która po jego zakończeniu jest automatycznie wycofywana (\textit{rollback}). Zapewnia to pełną izolację między testami bez konieczności ręcznego czyszczenia stanu bazy. Wzorzec ten zastosowano w testach konsumenta zdarzeń mikroserwisu \textit{tasks}:

\begin{lstlisting}[language=java, caption={Test konsumenta zdarzeń z izolacją @TestTransaction (apps/tasks/src/test/.../UserEventConsumerTest.java)}, label={lst:test-transaction}]
@QuarkusTest
class UserEventConsumerTest {

    @Inject
    UserEventConsumer consumer;

    @Test
    @TestTransaction
    void createdEventPersistsKnownUser() {
        consumer.consume(new UserEvent(99L, UserEvent.Type.CREATED));
        assertNotNull(KnownUser.findById(99L));
    }

    @Test
    @TestTransaction
    void deletedEventRemovesKnownUser() {
        KnownUser u = new KnownUser();
        u.id = 42L;
        KnownUser.persist(u);

        consumer.consume(new UserEvent(42L, UserEvent.Type.DELETED));
        assertNull(KnownUser.findById(42L));
    }
}
\end{lstlisting}

\subsubsection{Testy jako bramka w łańcuchu CI/CD}
Użycie testów w łańcuchu CI na początku jest celowe i ma bezpośredni wpływ na efektywność procesu (listing~\ref{lst:github-actions-build}). Krok \texttt{Run tests} poprzedza budowanie natywnego obrazu kontenera (\texttt{imageBuild}), które korzysta z kompilatora GraalVM AOT i trwa kilkanaście minut. Niepowodzenie testu natychmiast przerywa łańcuch -- bez zbędnego zużycia zasobów CI na kosztowną kompilację natywną.

\begin{itemize}
    \item \textbf{Zakres:} Testy pokrywają kluczowe komponenty domeny (logika REST, perzystencja przez Panache) oraz integrację z bazą danych uruchamianą przez Dev Services, bez konieczności mocków.
    \item \textbf{Anulowanie przestarzałych uruchomień:} Konfiguracja w sekcji \texttt{concurrency} z ustawieniem \texttt{cancel-in-progress: true} gwarantuje, że nowe uruchomienie łańcucha dla tej samej gałęzi automatycznie anuluje poprzednie, wciąż aktywne uruchomienie. Eliminuje to sytuacje, w których wyniki przestarzałych łańcuchów blokują lub dezorientują recenzentów kodu.
    \item \textbf{Izolacja per serwis:} Każdy mikroserwis posiada oddzielny plik workflow uruchamiany wyłącznie przy zmianach w jego katalogu. Testy serwisu \textit{users} nie są uruchamiane przy zmianie kodu \textit{tasks} -- co skraca czas informacji zwrotnej i obniża koszty CI.
\end{itemize}

\subsubsection{Widoczność wyników w procesie Pull Request}
GitHub Actions raportuje wyniki każdego zadania (\textit{job}) jako \textbf{status check} bezpośrednio w widoku Pull Requesta. Zadanie \texttt{Build, Test, Scan \& Push} dla każdego serwisu widoczne jest jako oddzielny, nazwany check -- programista natychmiast identyfikuje, który z mikroserwisów jest \textit{czerwony} i z jakiego powodu. Dopóki wszystkie aktywowane łańcuchy nie zakończą się sukcesem, scalenie PR jest zablokowane, co egzekwuje zasadę \textit{green build before merge} bez konieczności ręcznej interwencji recenzenta.
